\documentclass[aspectratio=169, 10pt]{beamer}

% ==========================================
% PAQUETES Y CONFIGURACIÓN
% ==========================================
\usepackage[utf8]{inputenc}
\usepackage[spanish, es-tabla]{babel}
\usepackage{amsmath, amssymb, bm}
\usepackage{booktabs}
\usepackage{tcolorbox}
\tcbuselibrary{skins, breakable}
\usepackage{listings}
\usepackage{tikz}
% Librería babel requerida para evitar colisiones con el símbolo '>' en español
\usetikzlibrary{positioning, arrows.meta, calc, babel}

% Tema y Colores (Estilo Académico Institucional)
\mode<presentation> {
    \usetheme{Madrid}
    % Azul oscuro institucional (tipo UCB)
    \definecolor{ucbblue}{RGB}{0, 51, 102} 
    \usecolortheme[named=ucbblue]{structure}
    \setbeamertemplate{navigation symbols}{} % Ocultar símbolos de navegación
    \setbeamertemplate{itemize items}[circle]
}

% Configuración para código XML (URDF)
\lstset{
    language=XML,
    basicstyle=\ttfamily\scriptsize,
    keywordstyle=\color{blue}\bfseries,
    stringstyle=\color{red!70!black},
    commentstyle=\color{green!60!black}\itshape,
    backgroundcolor=\color{gray!5},
    frame=single,
    rulecolor=\color{lightgray},
    showstringspaces=false
}

% ==========================================
% COMANDOS PERSONALIZADOS
% ==========================================
% 1. Entorno para Notas Pedagógicas del Docente
\newtcolorbox{notadocente}[1][Nota Pedagógica]{
    enhanced,
    colback=yellow!10!white,
    colframe=yellow!60!black,
    boxrule=0.5mm,
    arc=1mm,
    left=2mm, right=2mm, top=1mm, bottom=1mm,
    fonttitle=\bfseries\small,
    coltitle=black,
    title=\textbf{#1:}
}

% 2. Placeholder elegante para imágenes pendientes
\newcommand{\placeholderfig}[2][4cm]{
    \begin{center}
    \begin{tikzpicture}
        \node[draw=gray, dashed, thick, fill=gray!10, text=gray!70, 
              minimum width=0.9\linewidth, minimum height=#1, align=center, rounded corners] 
              {\footnotesize \textbf{ESPACIO PARA IMAGEN:}\\#2};
    \end{tikzpicture}
    \end{center}
}

% ==========================================
% METADATOS DE LA PRESENTACIÓN
% ==========================================
\title[Robótica Industrial (IMT-342)]{\LARGE Robótica Industrial (IMT-342)}
\subtitle{\textbf{Unidad I: Morfología, Geometría del Espacio y Cinemática Directa}}
\author[M.Sc. Bernardo Quiroga]{M.Sc. Ing. Bernardo Quiroga Turdera}
\institute[UCB Tarija]{Universidad Católica Boliviana ``San Pablo'' -- Sede Tarija}
\date{Semestre II-2026}

\begin{document}

% Portada
\begin{frame}
    \titlepage
\end{frame}

% ==========================================
% SEMANA 1 - CLASE 1
% ==========================================
\section{Morfología de Robots Industriales}

\begin{frame}{Introducción y Morfología}
    \framesubtitle{Arquitectura Física del ABB IRB 120}
    
    \begin{columns}[T]
        \begin{column}{0.5\textwidth}
            \textbf{Conceptos Clave:}
            \begin{itemize}
                \item Definición de robot manipulador articulado (cadena cinemática abierta).
                \item Estructura del semestre: Marco CDIO (Concebir, Diseñar, Implementar, Operar).
                \item El Proyecto Final Integrador (PFI) como eje central.
            \end{itemize}
        \end{column}
        \begin{column}{0.5\textwidth}
            \placeholderfig[3.5cm]{Esquema 3D del manipulador\\ABB IRB 120 indicando sus 6 ejes}
        \end{column}
    \end{columns}
    
    \vspace{0.3cm}
    \begin{notadocente}
        Aclarar a los estudiantes que cada ecuación teórica vista en la pantalla se programará en ROS y se ejecutará en el robot real del laboratorio.
    \end{notadocente}
\end{frame}

\begin{frame}{Eslabones, Articulaciones y Grados de Libertad}
    \framesubtitle{Criterio de Grübler-Kutzbach}
    
    \begin{columns}[T]
        \begin{column}{0.55\textwidth}
            \textbf{Anatomía de la Cadena Cinemática:}
            \begin{itemize}
                \item \textbf{Juntas 1, 2 y 3 (Estructura Principal):} Posicionamiento grueso del extremo en el espacio $\mathbb{R}^3$.
                \item \textbf{Juntas 4, 5 y 6 (Muñeca Esférica):} Orientación del efector final. Intersección en un punto común $p_w$.
            \end{itemize}
            
            \vspace{0.2cm}
            \textbf{Criterio de Movilidad espacial ($\lambda=6$):}
            \begin{equation*}
                \text{GDL} = \lambda (N - 1 - j) + \sum_{i=1}^{j} f_i
            \end{equation*}
        \end{column}
        \begin{column}{0.45\textwidth}
            \placeholderfig[4cm]{Diagrama mecánico explotado\\mostrando las 6 juntas (R)}
        \end{column}
    \end{columns}
    
    \begin{notadocente}[Nota para el Profesor]
        Dibujar en la pizarra el punto $p_w$ (intersección 4, 5 y 6). Explicar que este detalle de diseño es el que permitirá resolver la cinemática inversa analíticamente en la Semana 6.
    \end{notadocente}
\end{frame}

% ==========================================
% SEMANA 2 - CLASE 1
% ==========================================
\section{Geometría de Rotación y Grupo SO(3)}

\begin{frame}{El Grupo Especial Ortogonal $SO(3)$}
    \framesubtitle{Matrices de Rotación en $2\text{D}$ y $3\text{D}$}
    
    \begin{columns}[c]
        \begin{column}{0.4\textwidth}
            \placeholderfig[3cm]{Mapeo de vectores entre\\marco $\{B\}$ y marco $\{A\}$}
        \end{column}
        \begin{column}{0.6\textwidth}
            \textbf{Matriz de Rotación Directa:}
            \renewcommand{\arraystretch}{1.3} % Da respiro a la matriz
            \begin{equation*}
                ^{A}R_{B} = \begin{bmatrix} 
                \hat{x}_B \cdot \hat{x}_A & \hat{y}_B \cdot \hat{x}_A & \hat{z}_B \cdot \hat{x}_A \\ 
                \hat{x}_B \cdot \hat{y}_A & \hat{y}_B \cdot \hat{y}_A & \hat{z}_B \cdot \hat{y}_A \\ 
                \hat{x}_B \cdot \hat{z}_A & \hat{y}_B \cdot \hat{z}_A & \hat{z}_B \cdot \hat{z}_A 
                \end{bmatrix}
            \end{equation*}
        \end{column}
    \end{columns}
    
    \vspace{0.4cm}
    \textbf{Propiedades Fundamentales de $SO(3)$:}
    \begin{itemize}
        \item \textbf{Ortogonalidad:} $R^T R = I \implies R^{-1} = R^T$ (La inversa es la transpuesta).
        \item \textbf{Preservación de norma:} $\Vert R \cdot p \Vert = \Vert p \Vert$ (Transformación isométrica).
        \item \textbf{Determinante propio:} $\det(R) = +1$ (Conserva sistemas dextrógiros).
    \end{itemize}
\end{frame}

\begin{frame}{Ejes Fijos vs. Ejes Móviles}
    \framesubtitle{Composición de Rotaciones Sucesivas}
    
    \begin{columns}[T]
        \begin{column}{0.5\textwidth}
            \begin{block}{Rotación respecto a Ejes Fijos}
                (Sistema Base $\{A\}$)
                
                \vspace{0.2cm}
                \textbf{Regla:} Pre-multiplicación (por la izquierda).
                $$R_{total} = R_z \cdot R_y \cdot R_x$$
            \end{block}
        \end{column}
        \begin{column}{0.5\textwidth}
            \begin{block}{Rotación respecto a Ejes Móviles}
                (Sistema del Cuerpo $\{B\}$)
                
                \vspace{0.2cm}
                \textbf{Regla:} Post-multiplicación (por la derecha).
                $$R_{total} = R_x \cdot R_y \cdot R_z$$
            \end{block}
        \end{column}
    \end{columns}

    \vspace{0.2cm}
    \placeholderfig[2cm]{Comparativa de diagramas de árbol de multiplicación matricial}
    
    \begin{notadocente}[Nota para el Profesor]
        Hacer un ejercicio rápido en directo sosteniendo un objeto (ej. borrador). Girarlo respecto al aula (eje fijo) y luego respecto a su propia cara (eje móvil) para fijar el concepto visualmente.
    \end{notadocente}
\end{frame}

% ==========================================
% SEMANA 2 - CLASE 2
% ==========================================
\section{Parametrización Angular y Gimbal Lock}

\begin{frame}{Parametrización Mínima de la Orientación}
    \framesubtitle{Ángulos de Euler (Convención ZYZ)}
    
    \begin{columns}[T]
        \begin{column}{0.4\textwidth}
            \placeholderfig[3.5cm]{Secuencia de giros $(\phi, \theta, \psi)$\\sobre giroscopio 3D}
            
            \vspace{0.2cm}
            \textbf{Concepto Clave:}\\
            Representación de 9 elementos matriciales con solo \textbf{3 parámetros independientes}.
        \end{column}
        \begin{column}{0.6\textwidth}
            \textbf{Matriz Euler ZYZ:}
            
            \vspace{0.2cm}
            % Reducción de fuente para acomodar la matriz ancha
            {\footnotesize
            \renewcommand{\arraystretch}{1.5}
            $R_{ZYZ}(\phi, \theta, \psi) =$
            \begin{equation*}
                \begin{bmatrix} 
                c_\phi c_\theta c_\psi - s_\phi s_\psi & -c_\phi c_\theta s_\psi - s_\phi c_\psi & c_\phi s_\theta \\ 
                s_\phi c_\theta c_\psi + c_\phi s_\psi & -s_\phi c_\theta s_\psi + c_\phi c_\psi & s_\phi s_\theta \\ 
                -s_\theta c_\psi & s_\theta s_\psi & c_\theta 
                \end{bmatrix}
            \end{equation*}
            }
        \end{column}
    \end{columns}
\end{frame}

\begin{frame}{Pérdida de un GDL: Gimbal Lock}
    \framesubtitle{Singularidades en la Representación Angular}
    
    \begin{columns}[T]
        \begin{column}{0.4\textwidth}
            \placeholderfig[3cm]{Esquema de alineación del 1er y 3er eje (Ej: $\theta = 0$)}
        \end{column}
        \begin{column}{0.6\textwidth}
            \textbf{Análisis Matemático:}
            Si el segundo giro es nulo ($\theta = 0$), evaluamos la matriz:
            {\small
            \begin{equation*}
                R_{ZYZ} = \begin{bmatrix} 
                \cos(\phi + \psi) & -\sin(\phi + \psi) & 0 \\ 
                \sin(\phi + \psi) & \cos(\phi + \psi) & 0 \\ 
                0 & 0 & 1 
                \end{bmatrix}
            \end{equation*}
            }
        \end{column}
    \end{columns}

    \vspace{0.3cm}
    \begin{alertblock}{Consecuencia Crítica}
        Las variables $\phi$ y $\psi$ quedan acopladas matemáticamente en una única suma $(\phi + \psi)$. Es imposible distinguir giros independientes del primer y tercer eje. Esto justifica el salto hacia \textbf{Cuaterniones Unitarios} en el software de control.
    \end{alertblock}
\end{frame}

% ==========================================
% SEMANA 3 - CLASE 1
% ==========================================
\section{Operadores Homogéneos en SE(3)}

\begin{frame}{Matrices de Transformación Homogénea}
    \framesubtitle{El Grupo Especial Euclidiano $SE(3)$}
    
    \begin{columns}[c]
        \begin{column}{0.4\textwidth}
            \textbf{Composición Encadenada:}\\
            Para marcos múltiples, la transformación se acumula:
            $$^{A}T_{C} = {}^{A}T_{B} \cdot {}^{B}T_{C}$$
            
            \vspace{0.3cm}
            \textbf{Transformación de Puntos:}\\
            Usando coordenadas homogéneas:
            $$P_A = {}^{A}T_{B} \cdot P_B$$
        \end{column}
        \begin{column}{0.6\textwidth}
            \textbf{Estructura Analítica de $SE(3)$:}
            \renewcommand{\arraystretch}{1.2}
            \begin{equation*}
                ^{A}T_{B} = 
                \begin{bmatrix} 
                R_{3\times3} & p_{3\times1} \\ 
                0_{1\times3} & 1 
                \end{bmatrix} 
                = 
                \begin{bmatrix} 
                r_{11} & r_{12} & r_{13} & x \\ 
                r_{21} & r_{22} & r_{23} & y \\ 
                r_{31} & r_{32} & r_{33} & z \\ 
                0 & 0 & 0 & 1 
                \end{bmatrix}
            \end{equation*}
            \vspace{0.1cm}
            \placeholderfig[1.5cm]{Matriz articulada en bloques (R, p, Perspectiva, Escala)}
        \end{column}
    \end{columns}
\end{frame}

\begin{frame}{Inversión Eficiente sin Descomposición LU}
    \framesubtitle{Propiedades Matriciales Homogéneas}
    
    \textbf{Derivación Analítica de la Inversa:}
    Dada la matriz de transformación homogénea directa $^{A}T_{B}$, su inversa matemática exacta (sin usar algoritmos numéricos costosos) es:
    
    \begin{equation*}
        ({}^{A}T_{B})^{-1} = {}^{B}T_{A} = 
        \begin{bmatrix} 
        R^T & -R^T p \\ 
        0_{1\times3} & 1 
        \end{bmatrix}
    \end{equation*}
    
    \vspace{0.5cm}
    \begin{notadocente}[Nota para el Profesor]
        Demostrar en la pizarra que intentar calcular la inversa de una matriz de $4 \times 4$ mediante determinantes clásicos ($1/|A| \cdot \text{Adj}$) consume demasiada CPU. Usar la fórmula analítica con el bloque de traslación transpuesto $-R^T p$ reduce el cómputo en ROS a microsegundos.
    \end{notadocente}
\end{frame}

% ==========================================
% SEMANA 3 - CLASE 2
% ==========================================
\section{Cuaterniones e Interpolación SLERP}

\begin{frame}{Cuaterniones de Hamilton}
    \framesubtitle{Parametrización Global Sin Singularidades}
    
    \begin{columns}[T]
        \begin{column}{0.5\textwidth}
            \placeholderfig[3.5cm]{Giro de ángulo $\theta$ sobre\\un eje vectorial $\hat{k}$}
            \vspace{0.2cm}
            
            \textbf{Justificación ROS:}\\
            Estructura nativa (\texttt{geometry\_msgs/Quaternion}) para evitar el bloqueo de cardán en tiempo real.
        \end{column}
        \begin{column}{0.5\textwidth}
            \textbf{Definición del Cuaternión:}
            \begin{align*}
                \mathbf{q} &= [\eta, \boldsymbol{\epsilon}] \\
                &= \left[ \cos\left(\frac{\theta}{2}\right), \hat{k} \sin\left(\frac{\theta}{2}\right) \right] \\
                &= q_w + q_x \mathbf{i} + q_y \mathbf{j} + q_z \mathbf{k}
            \end{align*}
            
            \vspace{0.3cm}
            \textbf{Restricción de Unidad (Norma 1):}
            \begin{equation*}
                \Vert \mathbf{q} \Vert = \sqrt{q_w^2 + q_x^2 + q_y^2 + q_z^2} = 1
            \end{equation*}
        \end{column}
    \end{columns}
\end{frame}

\begin{frame}{Interpolación Esférica (SLERP)}
    \framesubtitle{Suavizado de Orientación en Trayectorias}
    
    \begin{columns}[T]
        \begin{column}{0.4\textwidth}
            \placeholderfig[3cm]{Trayectoria curva sobre una\\hiper-esfera unitaria $\mathbb{S}^3$}
        \end{column}
        \begin{column}{0.6\textwidth}
            \textbf{Algoritmo SLERP}\\
            Para conectar $\mathbf{q}_0$ y $\mathbf{q}_1$ en un tiempo normalizado $t \in [0,1]$:
            \vspace{0.1cm}
            {\small
            \begin{equation*}
                \text{SLERP}(\mathbf{q}_0, \mathbf{q}_1; t) = \frac{\sin((1-t)\Omega)}{\sin\Omega} \mathbf{q}_0 + \frac{\sin(t\Omega)}{\sin\Omega} \mathbf{q}_1
            \end{equation*}
            }
            \begin{equation*}
                \text{donde: } \cos\Omega = \mathbf{q}_0 \cdot \mathbf{q}_1
            \end{equation*}
            \vspace{0.1cm}
            \textbf{Concepto Clave:} Garantiza una \textit{velocidad angular constante} durante la reorientación de la pinza neumática.
        \end{column}
    \end{columns}
\end{frame}

% ==========================================
% SEMANA 4 - CLASE 1
% ==========================================
\section{Cinemática: Convención de Denavit-Hartenberg}

\begin{frame}{Protocolo Estándar de Denavit-Hartenberg (D-H)}
    \framesubtitle{Las 4 Reglas Geométricas de Asignación de Ejes}
    
    \begin{columns}[T]
        \begin{column}{0.45\textwidth}
            \placeholderfig[4.5cm]{Dos articulaciones consecutivas ($i$, $i+1$) con ejes locales}
        \end{column}
        \begin{column}{0.55\textwidth}
            \textbf{Reglas Absolutas de Construcción:}
            \begin{enumerate}
                \item \textbf{Eje $z_i$:} Se coloca a lo largo del eje de movimiento (revolución o prismático) de la junta $i+1$.
                \item \textbf{Eje $x_i$:} Se define a lo largo de la perpendicular común entre el eje $z_{i-1}$ y el eje $z_i$.
                \item \textbf{Eje $y_i$:} Se completa aplicando la regla de la mano derecha (producto cruz).
                \item \textbf{Origen $O_i$:} Se ubica en la intersección de la perpendicular común $x_i$ con el eje $z_i$.
            \end{enumerate}
        \end{column}
    \end{columns}
\end{frame}

\begin{frame}{Transformación Básica de Eslabón}
    \framesubtitle{Los 4 Parámetros Cinemáticos de Eslabón y Junta}
    
    \textbf{Ecuación General D-H ($^{i-1}T_{i}$):}
    \renewcommand{\arraystretch}{1.5}
    \begin{equation*}
        ^{i-1}T_{i} = 
        \begin{bmatrix} 
        \cos\theta_i & -\sin\theta_i \cos\alpha_i & \sin\theta_i \sin\alpha_i & a_i \cos\theta_i \\ 
        \sin\theta_i & \cos\theta_i \cos\alpha_i & -\cos\theta_i \sin\alpha_i & a_i \sin\theta_i \\ 
        0 & \sin\alpha_i & \cos\alpha_i & d_i \\ 
        0 & 0 & 0 & 1 
        \end{bmatrix}
    \end{equation*}

    \vspace{0.1cm}
    \placeholderfig[1.2cm]{Desglose de los giros y traslaciones $\theta_i, d_i, a_i, \alpha_i$}
    
    \begin{notadocente}[Mnemotecnia para Clase]
        Repetir con el grupo el orden exacto de los 4 movimientos base: Rotación $\theta_i$ en $z$, Traslación $d_i$ en $z$, Traslación $a_i$ en $x$, y Rotación $\alpha_i$ en $x$.
    \end{notadocente}
\end{frame}

% ==========================================
% SEMANA 4 - CLASE 2
% ==========================================
\section{Aplicación al Hardware: ABB IRB 120}

\begin{frame}{Diagrama Cinemático del ABB IRB 120}
    \framesubtitle{Plano Mecánico vs. Cero Matemático D-H}
    
    \begin{center}
        \placeholderfig[5.5cm]{Plano técnico oficial de ABB con cotas de longitud acoplado a ejes $z_0 \dots z_5$\\ ($c_1=290\text{mm}, a_1=270\text{mm}, c_3=302\text{mm}, c_4=72\text{mm}$)}
    \end{center}
\end{frame}

\begin{frame}{Parametrización del Manipulador del Laboratorio}
    \framesubtitle{Tabla de Parámetros D-H Definitiva}
    
    \begin{center}
        \renewcommand{\arraystretch}{1.3}
        \begin{tabular}{c c c c c}
            \toprule
            \textbf{Articulación $i$} & $\theta_i$ \textbf{(Variable)} & $d_i$ \textbf{(mm)} & $a_i$ \textbf{(mm)} & $\alpha_i$ \textbf{(rad)} \\
            \midrule
            1 & $q_1$ & 290 & 0 & $-\pi/2$ \\
            2 & $q_2 - \pi/2$ & 0 & 270 & 0 \\
            3 & $q_3$ & 0 & 70 & $-\pi/2$ \\
            4 & $q_4$ & 302 & 0 & $\pi/2$ \\
            5 & $q_5$ & 0 & 0 & $-\pi/2$ \\
            6 & $q_6$ & 72 & 0 & 0 \\
            \bottomrule
        \end{tabular}
    \end{center}
    
    \vspace{0.4cm}
    \begin{alertblock}{Concepto Crítico}
        Prestar especial atención a la justificación analítica del offset de $-\pi/2$ introducido en $\theta_2$. Esto se debe a la postura vertical erguida en la que se encuentra el robot físico cuando todos los encoders leen cero.
    \end{alertblock}
\end{frame}

% ==========================================
% SEMANA 5 - CLASE 1
% ==========================================
\section{Cinemática Directa y Modelado URDF}

\begin{frame}{El Modelo Cinemático Directo Completo}
    \framesubtitle{Multiplicación Encadenada de Eslabones}
    
    \begin{columns}[T]
        \begin{column}{0.4\textwidth}
            \placeholderfig[3.5cm]{Cadena de matrices acumuladas\\(Base $\to$ Brida)}
        \end{column}
        \begin{column}{0.6\textwidth}
            \textbf{Mapeo No Lineal ($^{0}T_{6}$):}
            
            \vspace{0.2cm}
            \begin{equation*}
                ^{0}T_{6}(\mathbf{q}) = {}^{0}T_{1}(q_1) \cdot {}^{1}T_{2}(q_2) \dots {}^{5}T_{6}(q_6)
            \end{equation*}
            donde $\mathbf{q} = [q_1, q_2, q_3, q_4, q_5, q_6]^T$.
            
            \vspace{0.4cm}
            \textbf{Concepto Clave:}\\
            Transformación determinista desde el \textit{espacio articular} $\mathbf{q} \in \mathbb{R}^6$ hacia el \textit{espacio cartesiano} de la tarea en $SE(3)$.
        \end{column}
    \end{columns}
\end{frame}

\begin{frame}[fragile]{Estructura de Marcado XML en Robótica}
    \framesubtitle{Anatomía de un Archivo URDF para ROS/RViz}
    
    \begin{columns}[T]
        \begin{column}{0.45\textwidth}
            \textbf{Estructura Jerárquica:}
            \begin{itemize}
                \item \texttt{<link>}: Eslabones estáticos (Masa, Inercia, STL).
                \item \texttt{<joint>}: Ejes de movimiento que conectan parent y child.
            \end{itemize}
            \vspace{0.2cm}
            \placeholderfig[1.5cm]{Árbol URDF visual}
        \end{column}
        \begin{column}{0.55\textwidth}
            \textbf{Fragmento Código URDF (.xml):}
\begin{lstlisting}
<joint name="joint_2" type="revolute">
  <parent link="link_1"/>
  <child link="link_2"/>
  <origin xyz="0 0 0.290" rpy="0 0 0"/>
  <axis xyz="0 1 0"/>
  <limit lower="-1.919" upper="1.919" 
         effort="100" velocity="3.14"/>
</joint>
\end{lstlisting}
        \end{column}
    \end{columns}

    \vspace{0.2cm}
    \begin{notadocente}[Precaución Laboratorio]
        La etiqueta \texttt{<origin>} en el URDF utiliza la transformación de origen fijo del CAD, mientras que la tabla D-H expresa la cinemática propia de eslabón. Conectar ambas representaciones para evitar confusión al visualizar en RViz.
    \end{notadocente}
\end{frame}

\end{document}